\section{Introduction}
Critic-free objectives such as RLOO and GRPO learn from terminal rewards
without a value model \citep{ahmadian2024back,shao2024deepseekmath}.
In language-model training, however, a response-level reward produces many
token updates. Response-level centering weights each response equally, whereas
expanded token advantages weight responses by their active lengths. A second issue
arises when rollout and actor policies differ: a per-token correction uses
only the current ratio, while a whole-response mask makes the same
decision for every position. These issues motivate REINFORCE Pro Max,
which combines sign-dependent token normalization with causal prefix masking.

We analyze reward-improvement conditions with sampled RLOO groups,
numerical safeguards, current-to-rollout ratios, and the active-token denominator.

\paragraph{Method and mathematical results.}
REINFORCE Max scales positive and negative advantages separately.
Its ideal factors balance signed token mass and set the second moment to one. We derive its binary-reward
coefficients and conditions under which the safeguarded update remains a
positive multiple of the success-probability gradient. Variable within-class
lengths and shared parameters across prompts introduce explicit covariance
terms. The ideal binary rule equals matched token-weighted standardization;
the safeguarded coefficients additionally depend on lengths and numerical
regularization.

REINFORCE Pro selects each position using the current-to-rollout prefix geometric mean
and retains a token-level importance coefficient. Its second-moment bound
separates an envelope times the accepted importance mass from a lower-side
reentry term. This term accounts for accepted tokens whose preceding prefix lies
below the lower threshold.

\paragraph{Relation to prior theory.}
PPO constrains token ratios through clipping \citep{schulman2017proximal}.
We instead use a masked importance-weighted objective and DPPO's terminal-reward decomposition
\citep{qi2026rethinkingtrustregionllm} and TRM's Adaptive and Mixed bounds
\citep{li2025trm} to connect this surrogate to reward.
CTPO~\citep{zhang2026ctpo} uses cumulative ratios; MinPRO~\citep{lei2026minpro}
clips and detaches the product of the current ratio and smallest preceding
ratio. PNPO~\citep{zhang2026pnpo}
uses current-to-behavior prefix geometric means as both weights and mask
statistics. CPPO~\citep{mao2026cppo} masks token contributions using
position-weighted divergence thresholds and cumulative prefix budgets.
REINFORCE Pro uses current-to-rollout geometric means only for
masking and retains a single token-level importance ratio.
These choices motivate our retained-coefficient and reward-error analyses.
The terminal RLOO signal depends on continuations and sibling responses;
we integrate its baseline over complete prompt groups. Prefix
change-of-measure identities apply directly to prefix-measurable functions;
our outcome signals require the group conditioning argument.
The normalization analysis of \citet{ge2026normalization} uses exact
on-policy gradients scaled by population reward standard deviations. Here
coefficients and token counts depend jointly on sampled responses.
Length-only weighting tradeoffs~\citep{ding2026length} give a related perspective.

For the combined method, an exact objective decomposition isolates the
scaling and masking costs. A fixed reference scale removes the
distortion of matching common scaling under full retention. Applying the
DPPO/TRM reward bounds gives a population improvement condition. In a
two-step family, this bound stays positive and retained second moments
remain bounded as unmasked moments grow without bound
(\Cref{ex:selective-positive-margin}). A shared-parameter counterexample
shows that an increasing masked objective can accompany decreasing reward.

\section{Finite model and the population reference}
\label{sec:submission-model}
Prompts follow a fixed finite law. A policy is a normalized kernel on a
finite token vocabulary and complete autoregressive histories
$c_t=(x,y_{<t})$. Responses terminate by a fixed cap $T$. After EOS, all
policies share an absorbing transition and the action mask is zero; the
EOS-emitting action itself remains active. Rewards satisfy
$|R(x,y)|\le R_{\max}$, and $J(q)=\mathbb E_{x,y\sim q}R(x,y)$.
The rollout policy $p=\piroll$ and evaluated policy $q=\pi_\theta$ have mutual
support at every context. At a represented token, define the single ratio
\[
 \rho_t(q)=\frac{q(y_t\mid c_t)}{p(y_t\mid c_t)}.
\]
Both the prefix statistic and the token coefficient use this ratio.
We analyze the masked surrogate without PPO clipping or a separate
importance-correction multiplier. Rewards are sparse and terminal;
KL, entropy, and auxiliary-loss coefficients are zero. The rollout law
and sampled advantages remain fixed during an update, whereas the mask
is recomputed from the current policy and detached during differentiation.

\paragraph{Raw surrogate and reward error.}
Let $b$ be independent of its own response conditional on the prompt, with
the rollout and baseline laws fixed during evaluation. The centered raw
surrogate is
\[
 \mathcal S_p(q)=\mathbb E_p\!\left[(R-b)\sum_t m_t(\rho_t-1)\right]
               =\mathbb E_p\!\left[R\sum_t m_t(\rho_t-1)\right].
\]
Conditional ratio normalization cancels the baseline, and the tower property
identifies this with the per-step advantage surrogate
(\Cref{prop:surrogate-equivalence}). Normalization, masking, and finite-batch
reduction enter separately through the objective decomposition below.

Write $\bar D_t=\mathbb E_{x,c_t\sim p}\operatorname{TV}(p(\cdot\mid c_t),
q(\cdot\mid c_t))$. Assume local TV and $D_{\rm KL}(p\|q)$ are bounded
by $\varepsilon$ and $\delta$ at every context before $T$, respectively.
Let $\bar D_{\rm seq}=\E_x\operatorname{TV}(P_x^p,P_x^q)$ for the complete
response laws. Combining the Adaptive and Mixed bounds of
\citet{li2025trm}, as proved in \Cref{thm:adaptive-bound}, gives
\begin{align*}
 |J(q)-J(p)-\mathcal S_p(q)|&\le B(p,q),\\
 B(p,q)&=\min\{B_{\rm Adap},B_{\rm Mix}\},\\
 B_{\rm Adap}&=4R_{\max}\sum_{t=1}^{T}\bar D_t
       \min\{1,(T-t)\varepsilon,\sqrt{(T-t)\delta/2}\},\\
 B_{\rm Mix}&=4R_{\max}T\min\{1,\varepsilon\}\bar D_{\rm seq}.
\end{align*}
Adaptive controls conditional continuation TV; Mixed controls marginal
prefix shifts. Both use policy-level bounds at all contexts before $T$,
distinct from sampled mask decisions. \Cref{sec:preliminaries} proves both
bounds and derives their DPPO and Pinsker--Marginal specializations.

\section{REINFORCE Max}
Max scales positive and negative advantages by separate positive factors
$\alpha$ and $\beta$. In this paper, its inputs are RLOO signals
$A_i=R_i-(n-1)^{-1}\sum_{j\ne i}R_j$ for $n\ge2$ responses to one prompt.
Sparse terminal reward with $\gamma=1$ assigns $A_i$ to every active token.
Let $P=\{(i,t):A_i>0,m_{i,t}=1\}$ and
$N=\{(i,t):A_i<0,m_{i,t}=1\}$. For nonempty $P,N$, we choose $\alpha,\beta>0$
as the ideal factors balancing signed token mass and setting the second moment to one:
\[
 \sum_{P}\alpha A_i+\sum_N\beta A_i=0,\qquad
 \frac{1}{|P|+|N|}\left(\sum_P(\alpha A_i)^2+\sum_N(\beta A_i)^2\right)=1.
\]
The normalized signal is $H_{i,t}=\alpha A_i$ on $P$, $\beta A_i$ on $N$,
and zero on raw zero tokens. Consequently signs and within-sign ordering are
preserved while positive and negative token mass are balanced. Numerical safeguards introduce the moment residuals in
\Cref{prop:normalization-safeguards}, retained in the objective decomposition.

\paragraph{Closed-form factors.}
For nonempty $P$ and $N$, let $S_+=\sum_P A_i>0$,
$S_-=-\sum_N A_i>0$, and $Q_\pm$ be the corresponding sums of squared
advantages. Writing $m=|P|+|N|$, the two constraints give
\[
 \kappa=\frac{S_+}{S_-},\qquad
 \alpha=\sqrt{\frac{m}{Q_++\kappa^2Q_-}},\qquad
 \beta=\kappa\alpha.
\]
Zero signed mass gives $\alpha S_+=\beta S_-$; the second-moment
constraint then determines the unique positive solution. Positive factors
preserve individual signs but can change the summed gradient direction.
For an exact RLOO group, $\sum_i A_i=0$, whereas the expanded
token sum is $\sum_i T_iA_i$. Unequal response lengths weight each
advantage differently; the asymmetric factors balance this token mass
before importance weighting and multiplication by score vectors.

\begin{proposition}[Binary RLOO normalization and response-length balance]
\label{prop:binary-rloo-normalization}
Suppose a prompt group contains $k\ge1$ responses with scalar reward $r_+$
and $m\ge1$ with reward $r_-<r_+$, and write $n=k+m$ and
$\Delta=r_+-r_->0$. Work in the RLOO branch and use the ideal factors in
Eq.~\eqref{eq:alpha-beta}. Let $P$ and $N$ be the total active-token counts
in the higher- and lower-reward responses, respectively; assume $P,N>0$.
Then the raw RLOO advantages are $a_+$ and $-a_-$, where
\[
    a_+=\frac{m\Delta}{n-1},\qquad
    a_-=\frac{k\Delta}{n-1}.
\]
Every active token in the two classes has normalized advantage
\begin{equation}
    \hat A_+=\alpha a_+=\sqrt{N/P},\qquad
    \hat A_-=-\beta a_-=-\sqrt{P/N}.
    \label{eq:binary-normalized-amplitudes}
\end{equation}
Writing $\bar T_+=P/k$ and $\bar T_-=N/m$ for the class-conditional mean
response lengths, the factor ratio and aggregate coefficient balance satisfy
\begin{equation}
    \frac{\beta}{\alpha}=\kappa
      =\frac{Pa_+}{Na_-}=\frac{\bar T_+}{\bar T_-},\qquad
    P\hat A_+=-N\hat A_-.
    \label{eq:binary-length-balance}
\end{equation}
These conclusions concern scalar advantages before importance weighting,
prefix masking, or multiplication by score vectors.
\end{proposition}
\begin{proof}[Proof of \Cref{prop:binary-rloo-normalization}]
Subtracting the sibling mean from each reward gives
$a_+=m\Delta/(n-1)>0$ and $a_-=k\Delta/(n-1)>0$.
The length-weighted first and second sums are $Pa_+,Na_-$ and
$Pa_+^2,Na_-^2$. Hence
\[
 \kappa=\frac{Pa_+}{Na_-},\qquad
 D=Pa_+^2+\kappa^2Na_-^2=a_+^2\frac{P(P+N)}{N}.
\]
Substituting into $\alpha=\sqrt{(P+N)/D}$ and $\beta=\alpha\kappa$
gives $\alpha a_+=\sqrt{N/P}$ and $\beta a_-=\sqrt{P/N}$.
Also $\kappa=Pm/(Nk)=(P/k)/(N/m)=\bar T_+/\bar T_-$, and
$P\sqrt{N/P}=N\sqrt{P/N}$ proves aggregate balance.
\end{proof}

\paragraph{Update direction under class-dependent lengths.}
Fix one prompt with rewards $R\in\{0,1\}$ and $n\ge2$ iid on-policy
responses. Let $s(\theta)=\Pr_{p_\theta}(R=1)$,
$g=\nabla s(\theta)$ and $u_\theta(y)=\nabla\log p_\theta(y)$.
Assume class-dependent lengths $0<L_0,L_1\le L_{\max}$, accepted prefix
masks. With the safeguarded factors of
\Cref{prop:binary-complete-update}, lower clamp $\eta\in(0,1]$, and group
token count $L$, the response-score coefficient is $C_i=H_i/L$.
The expected update satisfies
\[
 U:=\mathbb E\!\left[\sum_i C_i(Y_{1:n})u_\theta(Y_i)\right]
   =\Lambda g,\qquad \Lambda\ge\eta/L_{\max}>0.
\]
To see this, condition on the $n-1$ siblings $Z$ of response $i$.
The coefficient then depends on $Y_i$ only through its reward label.
The zero-mean score identity gives the multiplier
$\Lambda=\sum_i\mathbb E_Z[C_i(1,Z)-C_i(0,Z)]$.
The sibling mean lies in $[0,1]$, each positive scale is at least $\eta$,
and every group has at most $nL_{\max}$ active tokens. Each conditional
contrast is therefore at least $\eta/(nL_{\max})$; summing proves the bound.
The optional uniform replacement adds a nonnegative contrast on
all-success groups and zero on all-failure groups, preserving this bound.

Ideal normalization removes any common positive scaling of the raw
advantages. Hence mean-baseline and RLOO signals coincide after ideal
normalization; binary signals also equal matched token standardization
(\Cref{prop:normalization-scale-invariance,cor:binary-token-standardization}).

\paragraph{When the two scales change the update direction.}
The SGA analysis of \citet[Section~2]{li2025sga} separates gradient
alignment from stochastic noise. For Max, the relevant first-order
quantity is the inner product of the reward gradient with the expected
scaled update. The following decomposition identifies how dependence
between the scale and the response affects that alignment.

\begin{proposition}[Within-class length covariance decomposition]
\label{prop:binary-length-covariance}
Drop the class-dependent-length premise and let $C_i(y_{1:n})$ be the actual
response-score coefficient produced by the implementation, including its
length-weighted normalization, total-token denominator and optional
uniform-scale term. For a fixed direction $v$, write
$f_v(y)=\langle v,u_\theta(y)\rangle$. Assume both binary response classes
have positive finite probability. Define $\Lambda_C$ as the sum over
coordinates and sibling draws of the difference between the conditional
class means of $C_i$, and define
$\mathcal E_v$ as the corresponding weighted within-class covariance between
$C_i$ and $f_v$. Then the exact identity is
\begin{align*}
 \mathbb E\!\left[\sum_i C_i(Y_{1:n})f_v(Y_i)\right]
   &= \Lambda_C\,\mathbb E[\mathbf 1_{\{R=1\}}f_v(Y)] + \mathcal E_v,\\
 |\mathcal E_v|&\le\sqrt{V_C\,n\,V_{f,v}}.
\end{align*}
Here, writing $Z$ for the $n-1$ independent sibling responses and inserting
$Y$ at coordinate $i$, define
$\bar C_i(q,Z):=\mathbb E[C_i(Y,Z)\mid R(Y)=q,Z]$ and
$\bar f_v(q):=\mathbb E[f_v(Y)\mid R(Y)=q]$. The variances are
\begin{align*}
 V_C&:=\sum_i\mathbb E_{Z,Y}
       [(C_i(Y,Z)-\bar C_i(R(Y),Z))^2],\\
 V_{f,v}&:=\mathbb E_Y[(f_v(Y)-\bar f_v(R(Y)))^2].
\end{align*}
For $v=g=\nabla s(\theta)$, this gives
$\langle g,U\rangle\ge\Lambda_C\|g\|^2-\sqrt{V_C\,n\,V_{f,g}}$.
A positive right-hand side implies $s(\theta+tU)>s(\theta)$ for all
sufficiently small $t>0$ (\Cref{cor:binary-local-improvement}).
\end{proposition}

\begin{proof}
Condition on the sibling responses and split the current response support
into the two reward classes. Subtracting and adding the conditional class
means of both $C_i$ and $f_v$ gives the displayed identity; the terms with
one centered factor vanish within each class. Applying the finite
Cauchy--Schwarz inequality gives the square-root bound. The score identity
$\mathbb E[\mathbf 1_{\{R=1\}}f_g(Y)]=\|g\|^2$ yields the ascent bound;
differentiability gives local improvement. Class-dependent lengths make
$\mathcal E_v=0$, recovering \Cref{prop:binary-complete-update}.
\end{proof}
\paragraph{Multivalued rewards distinguish the two rules.}
\label{par:multivalued-sign-example}
For rewards $(0,2/3,1)$ and active lengths $(1,1,4)$, the actual RLOO
advantages are $(-5/6,1/6,2/3)$, with token mean $1/3$ and variance
$11/36$. Token standardization changes the second response's coefficient
from positive to negative, since $1/6-1/3=-1/6$; sign-dependent scaling
keeps $\alpha/6>0$. The sign change distinguishes token standardization from sign-dependent
scaling for multivalued rewards.

\paragraph{From a prompt direction to population ascent.}
For a fixed prompt with binary success probability $s_x(\theta)$, write
$g_x=\nabla s_x(\theta)$ and $U_x=\Lambda_xg_x+e_x$ for the conditional
population update. With the stated on-policy, accepted-mask
conditions, the covariance analysis bounds the projection of $e_x$
separately from the positive gradient multiplier $\Lambda_x$. Across a fixed finite prompt law,
let $G=\mathbb E_xg_x$, $U=\mathbb E_xU_x$, and
$f_x=\langle G,g_x\rangle$. Then \Cref{cor:prompt-ascent} gives
\[
 \langle G,U\rangle\ge
 \mathbb E_x\Lambda_x\|G\|^2-
 \sqrt{\operatorname{Var}_x(\Lambda_x)\operatorname{Var}_x(f_x)}
 -\mathbb E_x|\langle G,e_x\rangle|.
\]
A strictly positive right-hand side implies local reward improvement along
$U$ in the differentiable finite model. Differing positive prompt scales
can otherwise reweight conflicting gradients. The two subtracted terms
quantify the effects of between-prompt reweighting and within-class
coefficient variation, respectively.

\section{REINFORCE Pro: Causal Prefix Mask}
\subsection{Masked surrogate objective}
Following the masked objectives of DPPO and TRM
\citep{qi2026rethinkingtrustregionllm,li2025trm}, Pro uses
\[
 S_{\rm Pro}(q)=\sum_{i,t}\omega_{i,t}M_{i,t}(q)\rho_{i,t}(q)A_i,
 \qquad \rho_{i,t}(q)=\frac{q(y_{i,t}\mid c_{i,t})}{p(y_{i,t}\mid c_{i,t})}.
\]
Here $\omega_{i,t}=m_{i,t}/\sum_{j,s}m_{j,s}$ includes all active tokens,
including rejected ones. Pro Max replaces $A_i$ by the scaled signal $H_{i,t}$.

\paragraph{Prefix masking criterion.}
TRM assigns one mask to an entire response; Pro makes a separate decision
at each position using only its observed prefix. Define
\[
 g_{i,t}(q)=\exp\!\left(\frac{\sum_{s\le t}m_{i,s}\log\rho_{i,s}(q)}
 {\max(1,\sum_{s\le t}m_{i,s})}\right),\qquad
 M_{i,t}(q)=\mathbf1\{\ell_-\le g_{i,t}(q)\le\ell_+\}.
\]
For $0<\ell_-<1<\ell_+$, the geometric mean selects contributions;
$\rho_{i,t}$ supplies their token weight. Later prefixes may restore acceptance
without changing earlier decisions. DPPO instead uses a divergence- and
advantage-dependent mask. Signals and masks are detached during differentiation;
$\rho$ carries the gradient (Appendix~\ref{sec:unified}).

\paragraph{Exact objective decomposition.}
Let $k$ enumerate represented active positions, with nonnegative reduction
weights $\omega_k$, raw signal $A_k$, scaled signal $H_k$, and
$M_k=M_k(q)$, $\rho_k=\rho_k(q)$. Define
\begin{align*}
 U&=\sum_k\omega_k\rho_k A_k,&
 G&=\sum_k\omega_k\rho_k(1-M_k)A_k,\\
 N&=\sum_k\omega_k\rho_k M_k(H_k-A_k).
\end{align*}
\begin{proposition}[Finite-batch decomposition]
The masked objective satisfies
\[
 -\mathcal L=U-G+N,\qquad
 |-\mathcal L-U|\le |G|+|N|.
\]
The identity holds for every realized batch, including random response
lengths and policy-dependent masks.
\end{proposition}
Here $G$ and $N$ record rejection and scaling. Adding and subtracting
the retained raw contribution proves the identity
(\Cref{prop:promax-objective-split}). This adapts the discarded-contribution
analysis \citep[Section~1.3]{li2025sequencemasking} to a tokenwise prefix mask.

\subsection{A bound on retained-weight moments}
The quantity to control is the retained coefficient $M_t\rho_t$.
Individual token ratios can be arbitrarily large under prefix-average acceptance.
For a normalized two-action, two-step policy pair, the ratios on one path
can be $(K^{-1},K)$, so $M_2=1$ for any thresholds containing one, yet
\[
 \mathbb E_{\piroll}[(M_2\rho_2)^2]
 \ge \frac{K^3}{(K+1)^2}\ge K/4\quad(K\ge1).
\]
\Cref{prop:prefix-weight-limit} gives the construction. Controlling this
moment requires accounting for reentry into the acceptance interval.
The mask bounds a prefix average rather than each token ratio.
The reward bound assumes local TV and KL conditions on the policy pair.

\paragraph{Second-moment bound with lower-side reentry.}
For this bound, the rollout and current kernels have mutual support,
and $M_t$ denotes the mask on the full padded tree prefix.
Let $W_t=\prod_{s\le t}\rho_s$, $W_0=1$, $\tau_+=\log\ell_+$ and
$\tau_-=-\log\ell_-$. The lower-excursion contribution is
$\mathcal R_t^-=\mathbb E_{\piroll}[
\mathbf 1\{\log W_{t-1}<-(t-1)\tau_-\}M_t\rho_t^2]$.
If $\log W_{t-1}\ge-(t-1)\tau_-$ and the current token is accepted,
then $\rho_t\le C_t=e^{t\tau_++(t-1)\tau_-}$. Hence
\[
 \mathbb E_{\piroll}[(M_t\rho_t)^2]
 \le C_t\mathbb E_{\piroll}[M_t\rho_t]+\mathcal R_t^-
 \le C_t+\mathcal R_t^-.
\]
\emph{Proof, step 1: a local envelope.} Acceptance gives
$\log W_t\le t\tau_+$; subtracting the previous log-product bounds $\rho_t$.
\emph{Step 2: conditional normalization.} At a history $h$ above the
lower threshold, we have
\[
 \mathbb E_{\piroll}[M_t\rho_t^2\mid h]
 \le C_t\mathbb E_{\piroll}[M_t\rho_t\mid h]
 \le C_t\sum_yq(y\mid h)=C_t.
\]
\emph{Step 3: split the histories.} Controlled histories contribute at most
$C_t\mathbb E_{\piroll}[M_t\rho_t]$; the remaining histories contribute
$\mathcal R_t^-$. Normalization bounds the first moment by one.
Here $\mathbb E_{\piroll}[M_t\rho_t]$ is importance-weighted retained mass.
The argument adapts the Seq-MIS envelope analysis
\citep[Section~1.3]{li2025sequencemasking} to the current-policy prefix mask.

For $T\tau_+\le a$, $T\tau_-\le b$, $C_t\le e^{a+b}$ for $t\le T$.
Upper-side reentry gives $\rho_t\le e^{\tau_+}$; lower-side reentry
contributes the additional moment term. The same bound controls the active-token
coefficient $m_tM_t\rho_t$ (\Cref{thm:prefix-reentry-moment}).

\Cref{ex:prefix-half-retention} gives a nontrivial regime: the mask retains
half the rollout mass at every depth with $\mathcal R_t^-=0$ and second
moment $1/2$, while the unmasked first-token moment is arbitrarily large.

\paragraph{Causal rejection.}
The prefix log-average updates by a convex combination, allowing later
ratios to restore acceptance (\Cref{lem:prefix-stability}). Early mismatch
can suppress later terms; a deviation at $t^*$ leaves earlier decisions
intact. \Cref{thm:prefix-tighter} quantifies the rejection duration.

Algorithm~\ref{alg:promax} uses per-prompt Max scales and an active-token mean.
Token-count weighting preserves this objective; the implementation's
response-count weighting is treated in \Cref{app:promax-support}.

\begin{algorithm}[t]
\caption{REINFORCE Pro Max: core transformations.}
\label{alg:promax}
\begin{algorithmic}[1]
\REQUIRE Batch $(y_i,A_i,m_{i,t})$ with prompt-group index $b(i)$; policies $p,q$;
Pro interval $[\ell_-,\ell_+]$, step size $\eta$
\STATE $(\alpha_b,\beta_b)\leftarrow\operatorname{MaxScales}(A_b,m_b)$ for each prompt group $b$
\STATE $H_{i,t}\leftarrow\operatorname{stopgrad}\!\left(A_i(\alpha_{b(i)}\mathbf1_{A_i>0}+\beta_{b(i)}\mathbf1_{A_i<0})\right)$ \COMMENT{Max}
\STATE $z_{i,t}\leftarrow\log q(y_{i,t}\mid c_{i,t})-\log p(y_{i,t}\mid c_{i,t})$; $\rho\leftarrow\exp z$
\STATE $g_{i,t}\leftarrow\exp\!\left(\dfrac{\sum_{s\le t}m_{i,s}z_{i,s}}{\max(1,\sum_{s\le t}m_{i,s})}\right)$
\STATE $M\leftarrow\operatorname{stopgrad}(\mathbf1_{\ell_-\le g\le\ell_+})$ \COMMENT{Pro}
\STATE $\mathcal L\leftarrow-\dfrac{\sum_{i,t}m_{i,t}M_{i,t}\rho_{i,t}H_{i,t}}{\sum_{i,t}m_{i,t}}$
\STATE $\theta\leftarrow\theta-\eta\nabla_\theta\mathcal L$
\end{algorithmic}
\end{algorithm}

$\operatorname{MaxScales}$ uses the two factors derived above, with the
numerical safeguards in \Cref{prop:normalization-safeguards}; a fallback
returns $(1,1)$. The optional homogeneous-group replacement in
\Cref{app:uniform-scale} bypasses this core transformation.

\section{From the masked objective to reward improvement}
Following DPPO and TRM, we bound the raw reward error, quantify the
scaling and masking cost, and combine both in a reward-improvement bound.
\subsection{Reward decomposition and sequence error}
For a fixed candidate $q$ and rollout policy $p$, first isolate the
difference between true reward and the raw surrogate. We use DPPO's product decomposition
\citep[Appendix~B.1--B.2]{qi2026rethinkingtrustregionllm} and TRM's
factor-by-factor Adaptive proof \citep[Appendix~D]{li2025trm}.
\emph{Step 1: exact reward decomposition.} At a fixed prompt, telescoping gives
\[
 J_x(q)-J_x(p)-\mathcal S_{p,x}(q)
 =\E_p\!\left[R\sum_t(\rho_t-1)
                      \left(\prod_{j>t}\rho_j-1\right)\right].
\]
$J_x$ and $\mathcal S_{p,x}$ are prompt-conditional reward and raw surrogate.
\emph{Step 2: bound the continuation factor.} Conditioning on
$c_{t+1}=(x,y_{\le t})$ fixes the token factor and gives continuation TV:
\[
 \E_p\!\left[\left|\prod_{j>t}\rho_j-1\right|\mid c_{t+1}\right]
 =2\operatorname{TV}(P^p_{>t},P^q_{>t}).
\]
Both continuation laws condition on that same prefix. The sequence-TV
inequality bounds this quantity through subsequent local divergences:
\[
 \operatorname{TV}(P^p_{>t},P^q_{>t})
 \le\sum_{j>t}\E_{p(\cdot\mid c_{t+1})}
       \operatorname{TV}(p(\cdot\mid c_j),q(\cdot\mid c_j))
 \le(T-t)\varepsilon.
\]
To obtain the first inequality, telescope the two continuation probability
products and sum absolute differences. Each candidate suffix sums to one,
leaving rollout-weighted local TVs. Combining this with the unit TV bound
and KL chain rule followed by Pinsker gives
$b_{T-t}=\min\{1,(T-t)\varepsilon,\sqrt{(T-t)\delta/2}\}$.
\emph{Step 3: aggregate the error.} The tower property and
$\E_{x,y\sim p}|\rho_t-1|=2\bar D_t$ give $B_{\rm Adap}$.

\emph{Mixed bound: control the prefix distribution instead.} Following
\citet[Appendix~C.2]{li2025trm}, let $d_t^p,d_t^q$ be the prefix laws at a
fixed prompt and $g_t(c)=\E_{a\sim q(\cdot\mid c)}A_t^p(c,a)$, where
$A_t^p(c,a)=\E_p[R\mid c,a]-\E_p[R\mid c]$.
The performance-difference identity gives
\[
 J_x(q)-J_x(p)-\mathcal S_{p,x}(q)
 =\sum_t\bigl(\E_{d_t^q}g_t-\E_{d_t^p}g_t\bigr).
\]
Conditional normalization gives $\|g_t\|_\infty\le2R_{\max}\min(1,\varepsilon)$.
Marginalization gives $\operatorname{TV}(d_t^p,d_t^q)\le
\operatorname{TV}(P_x^p,P_x^q)$. Applying
$|\E_Pg-\E_Qg|\le2\|g\|_\infty\operatorname{TV}(P,Q)$, then summing and
averaging over prompts, gives $B_{\rm Mix}$. Both routes bound the same
error, so their minimum yields $B(p,q)$ (\Cref{app:adaptive-proof}).

\subsection{Scaling and masking error}
The rollout law, sampled RLOO signals and Max factors are fixed when
comparing the current policy $q$ with $p$. The mask is evaluated at each
policy: $M_t(p)=1$, since all reference ratios equal one.
For one group, write $L=\sum_{i,t}m_{i,t}>0$, $S_{\rm PM}=-\mathcal L$,
and $\Delta S_{\rm PM}=S_{\rm PM}(q)-S_{\rm PM}(p)$.
Group objectives enter the batch objective with weights proportional
to their active-token counts. For a fixed reference scale $\lambda>0$, define
\[
 D_\lambda(q)=\frac1L\sum_{i,t}m_{i,t}
 \left[|\rho_{i,t}-1|\,|\lambda A_i-H_{i,t}|
       +\rho_{i,t}(1-M_{i,t})|H_{i,t}|\right].
\]
The first term measures scaling discrepancy along the policy increment;
the second charges contributions rejected by the current-policy mask.
The scale $\lambda$ fixes comparison units and does not change the algorithm.
In particular, common scaling $H=\lambda A$ removes the first term,
and full retention removes the second. Both vanish at $q=p$.
Detaching $M$ during backpropagation does not make it constant between
candidate policies; this comparison includes its change explicitly.

\subsection{Population reward bound}
\begin{theorem}[Population reward improvement]
\label{thm:promax-population}
Assume the finite autoregressive model and the support and divergence
conditions of \Cref{thm:adaptive-bound}. For each prompt draw $n\ge2$
conditionally iid rollout responses and form their RLOO signals.
Use the same sampled advantages, scaling factors and active-token
reduction at $p$ and $q$, and compute each mask from its own policy ratio.
With expectation over the prompt and its response group, for any fixed
$\lambda>0$ let
\[
 \mathfrak M_\lambda(q)=\frac1{\lambda n}\E\!\left[
 L\bigl(\Delta S_{\rm PM}-D_\lambda(q)\bigr)\right].
\]
Then
\begin{equation}
 J(q)-J(p)\ge\mathfrak M_\lambda(q)-B(p,q).
 \label{eq:promax-population-margin}
\end{equation}
Thus $\mathfrak M_\lambda(q)>B(p,q)$ suffices for strict reward improvement.
\end{theorem}

\begin{proof}
\emph{Step 1: recover the raw increment.}
Let $Z=\sum_i A_i\sum_t m_{i,t}(\rho_{i,t}-1)$.
Since $M(p)=1$, the pointwise residual is
\begin{align*}
 \lambda Z-L\Delta S_{\rm PM}
 &=\sum_{i,t}m_{i,t}\bigl[(\rho_{i,t}-1)(\lambda A_i-H_{i,t})\\
 &\hspace{28mm}+\rho_{i,t}(1-M_{i,t})H_{i,t}\bigr].
\end{align*}
Its absolute value is at most $LD_\lambda(q)$, so
$L(\Delta S_{\rm PM}-D_\lambda(q))/\lambda\le Z$ on every group.

\emph{Step 2: integrate the group signal.} Conditional ratio normalization gives
$\E_p[\sum_t m_t(\rho_t-1)\mid x]=0$. Each sibling baseline is
independent of its own response given $x$, so
\[
 \frac1n\E[Z\mid x]
 =\E_{y\sim p(\cdot\mid x)}\!\left[R(x,y)\sum_t m_t(\rho_t-1)\right].
\]
Averaging over prompts identifies $\E[Z]/n=\mathcal S_p(q)$.
\emph{Step 3: subtract the sequence error.} The combined bound yields
$J(q)-J(p)\ge\mathcal S_p(q)-B(p,q)\ge\mathfrak M_\lambda(q)-B(p,q)$.
\end{proof}

For $K\ge1$ iid prompt groups, let $L_\Sigma=\sum_b L_b$ and, for
$F\in\{\Delta S_{\rm PM},D_\lambda\}$, define
$F_{\rm batch}=\sum_b L_bF_b/L_\Sigma$. Then the same margin is
\[
 \mathfrak M_\lambda(q)=\frac1{\lambda Kn}\E\!\left[
 L_\Sigma\bigl(\Delta S_{{\rm PM},{\rm batch}}-D_{\lambda,{\rm batch}}\bigr)\right].
\]
Multiplying by $L_\Sigma$ restores the group contributions; iid averaging
cancels $K$. The random denominator stays inside the expectation.
If $M=1$ and $H=\lambda A$, then $D_\lambda=0$.
Write $\mathfrak M=\mathfrak M_1$. The examples specify their reference
scale and use the full rollout law, including rejected branches.

\subsection{A positive regime and a failure boundary}
\begin{example}[Positive margin with nontrivial prefix rejection]
\label{ex:selective-positive-margin}
Fix $a=1-\zeta$, $0<\zeta<1/5$, $\tau=\log(5/4)$ and
$B>e^{2\tau}$. On three tokens, let
\[
 p_1=\left(a,\frac{\zeta}{B+1},\frac{\zeta B}{B+1}\right),\qquad
 q_{\delta,1}=\left(a,\frac{\zeta B}{B+1},\frac{\zeta}{B+1}\right).
\]
The rollout continuation is uniform. The current policy is uniform after
tokens $1,2$, and after token $0$ it is
$q_\delta(\cdot\mid0)=(1/3+\delta,1/3-\delta,1/3)$,
where $0\le\delta\le1/15$. Use mask interval $[4/5,5/4]$.
The first ratios are $(1,B,B^{-1})$; on branch $0$, second ratios are
$(1+3\delta,1-3\delta,1)$. The mask therefore accepts exactly branch $0$
at both positions. Set $R(y)=\mathbf1\{y_1=y_2=0\}$.
Then $J(q_\delta)-J(p)=a\delta$.

Use two iid rollout responses per group, four-token reduction $L=4$,
no group masking or \texttt{uniform\_scale}, stabilizer
$0\le\epsilon\le1/2$, product cap $10^8$, and factor clamp $[\epsilon,10]$.
Mixed-reward groups have $H=s_\epsilon A$, where
$s_\epsilon=\sqrt{4/(4+\epsilon)}\in[3/4,1]$;
equal-reward groups have zero signals. All rejected tokens remain in $L$.
For $f_q(y)=\sum_tM_t(q)\rho_t(q)$, independence of the sibling response gives
\[
 \E[S_{\rm PM}(q)]=\frac{s_\epsilon}{2}\operatorname{Cov}_p(R,f_q),
 \qquad S_{\rm PM}(p)=0.
\]
Here $\E_pR=a/3$, $\E_pf_q=2a$, and
$\E_p[Rf_q]=a(2+3\delta)/3$. Consequently,
\[
 \E[\Delta S_{\rm PM}]=s_\epsilon(a\zeta/3+a\delta/2),\qquad
 \E[D_{s_\epsilon}]=s_\epsilon a\zeta/3,
 \qquad \mathfrak M_{s_\epsilon}=a\delta.
\]
For the distortion, the scaling term vanishes. A rejected response has
reward zero, expected absolute RLOO signal $a/3$, and total rejected
ratio mass $2\zeta$, which gives the displayed expectation.

The root TV is $\bar D_1=\zeta(B-1)/(B+1)<\zeta$;
continuation TVs are at most $\delta$. For $\zeta\le\delta$, all local
TVs are at most $\delta$, so the Adaptive bound gives
$B_{\rm Adap}\le4\bar D_1\delta$.
Thus $\zeta=1/100$ and $\delta=1/30$ yield, uniformly in $B>25/16$,
\[
 J(q_\delta)-J(p)\ge\mathfrak M_{s_\epsilon}-B(p,q)
 \ge(1-5\zeta)\delta=\frac{19}{600}>0.
\]
At depths one and two the retained second moments are $a$ and
$a(1+6\delta^2)$, with no lower-side reentry. They remain bounded as
$\E_p[\rho_1^2]=a+\zeta(B^2+B^{-1})/(B+1)$ grows without bound.
For a parameter update within this mask region,
$\partial_\delta\E S_{\rm PM}=s_\epsilon a/2$ and
$\partial_\delta J=a>0$: the population ascent direction increases reward.
\end{example}

\begin{proposition}[Objective ascent can lower reward under a shared parameter]
\label{prop:coupled-gate-failure}
There is a normalized, strictly positive two-step current-policy family
with one shared parameter and a fixed rollout law such that arbitrarily
small positive parameter changes increase the expected ProMax objective
while decreasing reward. The mask has nontrivial prefix rejection.
\end{proposition}
\begin{proof}
Let $p_1=(1/2,1/6,1/3)$ with uniform continuations, and let
$q_{\delta,1}=(1/2,1/3,1/6)$. Set
\[
 q_\delta(\cdot\mid0)=(1/3+\delta,1/3-\delta,1/3),\qquad
 q_\delta(\cdot\mid1)=(1/3-3\delta,1/3+3\delta,1/3),
\]
with uniform continuation after token $2$ and $0\le\delta\le1/45$.
Use $[4/5,5/4]$ for the prefix mask and $R(y)=\mathbf1\{y_2=0\}$.
At the first token, the ratios are $(1,2,1/2)$. On branch $0$, the
second prefix products lie in $[14/15,16/15]$, inside the squared mask
interval. On branch $1$ they are at least $2(1-9\delta)\ge8/5>25/16$;
on branch $2$ they equal $1/2<16/25$. Thus both depths accept exactly
branch $0$ throughout this parameter interval.

With the same two-response groups and safeguards as above, $H=s_\epsilon A$.
The full rollout law remains fixed. The accepted ratio sum changes by
$3\delta\mathbf1_{y_1=0}(\mathbf1_{y_2=0}-\mathbf1_{y_2=1})$.
Its mean is zero and its covariance with $R$ is $\delta/2$. Therefore
\begin{equation}
 \E[S_{\rm PM}(q_\delta)-S_{\rm PM}(q_0)]
 =s_\epsilon\delta/4>0,\qquad
 J(q_\delta)-J(q_0)=-\delta/2<0.
 \label{eq:coupled-gate-failure}
\end{equation}
Here $q_0$ is the zero-parameter member of the same current-policy family,
not an additional importance-sampling reference. The negative reward
change follows from $\tfrac12\delta-\tfrac13(3\delta)$.
The rejected branch carries the larger adverse change. This example
separates local objective ascent from the sufficient population margin,
which also accounts for discarded contributions and sequence error.
\end{proof}

\section{Limitations}
The guarantees apply to finite autoregressive models with bounded terminal
rewards and the stated support, divergence, and reduction conditions.
Local ascent concerns on-policy population directions and requires positive
class masses and the covariance margins. The population reward bound
compares the current and rollout policies with sampled signals held fixed;
it includes the change of the prefix mask and the distortion from scaling. Its
applicability depends on the resulting margin relative to the sequence-error bound.

\section{Conclusion}
REINFORCE Pro Max combines sign-dependent RLOO scaling with causal prefix
masking. Its analysis connects sample-dependent scales to update alignment
and retained importance mass to a second-moment bound with lower-side reentry.
The reward bound restores the token normalization and subtracts scaling
and rejection costs from the population objective increment.
The smaller Adaptive or Mixed remainder then bounds the residual reward error. A two-step family has bounded retained second moments, a positive
reward bound, and arbitrarily large unmasked second moments. The shared-parameter
counterexample shows that objective ascent alone can lower reward. Both
the mask and the retained token weight compare the current policy directly with rollout.
